\documentclass[12pt, a4paper]{article}
\usepackage[T2A]{fontenc}
\usepackage{amsmath, amsfonts}
\usepackage[russian]{babel}
\begin{document}
Может ли число, составленное только из цифр 2 и 0, быть степенью выше первой натурального числа?

Сразу оговоримся, что 0 натуральным числом считать не будем~--- иначе задача будет тривиальной и неинтересной.

Пусть $q = a^p$, где $a,p,q\in \mathbb{N}$, $p > 1$, $a > 0$, число $q$ составленно только из 2 и 0.
Зафиксируем $p$ и возьмём минимально возможное $q$, удовлетворяющее этому соотношению.

Очевидно, что $q$ чётное, следовательно, $a$ тоже чётное.
А из этого в свою очередь следует что $q$ делится на $2^p$; в частности, $q$ делится на 4.
Согласно признаку делимости на 4, число составленное из двух последних цифр числа $q$, делится на 4.

Такими числами могут быть 00, 02, 20, 22. Очевидно, что нас устраивают только варианты $00=0$ и 20.
Заметим, что оба этих варианта заканчиваются на 0, а значит кратны 5.
Отсюда следует, что $q$ кратно 5, $a$ кратно 5, и наконец, $q$ кратно $5^p$.

Поскольку $a$ кратно 2 и 5, $a$ кратно 10 и заканчивается на один или более нулей.
Из тех же самых соображений $q$ кратно $10^p$ и заканчивается не меньше чем на $p$ нулей.
Иными словами, $a=10a_1$, $q=10^pq_1$, при этом $q_1$ составлено только из 2 и 0
(деление $q$ на степени десятки просто обрезает нули справа, другие цифры остаются теми же самыми).

В конечном итоге $q_1=a_1^p$, $0 < q_1 < q$, $q_1$ составлено только из 2 и 0;
а это явно противоречит нашей гипотезе, что $q$ \textit{минимальное} число с подобными свойствами.

Следовательно, такого $q$ не существует.
\end{document}
